\documentclass[12pt]{article}

% fonts

\usepackage[T1]{fontenc}
\usepackage{newtxtext}
\usepackage{bm}
\usepackage{amsthm} % before newtxmath b/c latex
\usepackage{newtxmath}
% \usepackage[bb=mt,cal=mt]{mathalpha}

\usepackage{amsmath}
\usepackage{amsfonts}
\usepackage{thmtools}
\usepackage{mathtools}

% spacing

\usepackage[parfill]{parskip}
\setlength{\parindent}{0pt}
\setlength{\parskip}{2ex plus 0.2ex minus 0.2ex}

% geometry of the page

\usepackage[
paper  = letterpaper,
left   = 1in,
right  = 1in,
top    = 1in,
bottom = 1in,
]{geometry}

% references

\usepackage{natbib}


\begin{document}

\begin{flushleft}
\textbf{Black Box Variational Inference (Ranganath et al., 2014)} \\
Satvik Saha (ss7228) \\
\today
\end{flushleft}

\vspace{0.1in}

``Black box''\footnote{%
    Prof.\ Blei has suggested in his book that ``turnkey'' better captures the
    intent: not mysterious, just mechanical.
    The rise of tooling that hides the details from practitioners may explain
    why ``black box'' stuck.
} variational inference solves the problem of optimizing the Evidence Lower
Bound (ELBO) $\mathcal{L}$ (over a variational family $\{q_\lambda\}$) by
`simply' performing (stochastic) gradient descent on the parameter $\lambda$.
This paper shows that VI is possible with access only to the model joint
density $p(x, z)$, variational densities $q_\lambda$ (with scores
$\nabla_\lambda \log{q_\lambda}$), and samples from $q_\lambda$: these combine
to give a Monte-Carlo estimate of $\nabla_\lambda \mathcal{L}$, which informs
the direction in which $\lambda$ moves.
A couple of variance reduction techniques (Rao-Blackwellization and control
variates) are also suggested to counter the noisiness of the estimate of
$\nabla_\lambda \mathcal{L}$.

In hindsight, it may seem surprising that such a direct approach (optimize
$\mathcal{L}(\lambda)$ directly via gradient descent) took time to gain
traction.
I would guess that Monte-Carlo estimation of $\nabla_\lambda\mathcal{L}$ was
seen as too noisy to be practical at the time.
The variance reduction ideas in this paper are therefore crucial to its
viability; the authors point out that earlier related work had struggled with
this issue.
Another factor may have been the field's emphasis on analytical, model-specific
updates (as in CAVI), in contrast with BBVI's generic stochastic optimization
approach.
I would be curious to learn more about how this paper was initially received,
given that today such methods feel almost indispensable.

% \nocite{*}
% \bibliographystyle{plainnat}
% \bibliography{week02}


\end{document}

%%% Local Variables:
%%% mode: latex
%%% TeX-master: t
%%% End:
